
\SGARefTarget{sga3:xi:section:6}
\section*{6. A Rigidity Property for Homomorphisms\\
of Certain Group Schemes, and the Representability\\
of Certain Transporters\footnote{\SGARefLink{sga3:xi:section:6}{The
present section} does not use the results of Sections \SGARefLink{sga3:xi:section:3}{3}, \SGARefLink{sga3:xi:section:4}{4}, and \SGARefLink{sga3:xi:section:5}{5}; its
natural place would be in Exposé VI\textsubscript{B}.}}
\addcontentsline{toc}{section}{6. A Rigidity Property for Homomorphisms
of Certain Group Schemes, and the Representability of Certain
Transporters}

\medskip
\noindent\SGARefTarget{sga3:xi:proposition:6:1}\textbf{Proposition 6.1.}
---
\textit{Let \(S\) be a locally noetherian prescheme, let \(H\) be a
commutative \(S\)-group prescheme of finite type over \(S\), and let
\(E\) be a set of positive integers stable under multiplication.
Suppose that the following conditions hold:}

\begin{enumerate}
\item[(a)] \SGARefTarget{sga3:xi:proposition:6:1:item:a}%
\textit{for every \(n\in E\), the subgroup}
\SGARefTarget{sga3:xi:display:6:1:n-torsion}
\[
  {}_nH=\operatorname{Ker}\bigl(n\cdot\operatorname{id}_H\bigr)
\]
\textit{is finite and flat over \(S\);}

\item[(b)] \SGARefTarget{sga3:xi:proposition:6:1:item:b}%
\textit{for every \(s\in S\), the family of the
\({}_nH_s\), \(n\in E\), is schematically dense in the fibre
\(H_s\).\footnote{The printed French text has \(H\) here; the
fibrewise condition and its use in the proof require \(H_s\).}}
\end{enumerate}

\textit{Let \(Y\) be a closed subprescheme of \(H\).  Then the
subfunctor \(\prod_{H/S}(Y/H)\) of \(S\) (compare \SGARefLink{sga3:viii:section:6}{Exposé VIII, 6}) is
representable by a closed subprescheme \(T\) of \(S\).  If \(S\) is
noetherian, there exists \(n\in E\) such that\footnote{The printed
formula writes the right-hand quotient as
\((Y\cap{}_nH)/S\).  Since the product functor is taken over
\({}_nH/S\), and as the proof requires, the quotient is
\((Y\cap{}_nH)/{}_nH\).}}
\SGARefTarget{sga3:xi:display:6:1:cross}
\[
  \tag{x}
  \prod_{H/S}(Y/H)
  =
  \prod_{{}_nH/S}\bigl((Y\cap{}_nH)/{}_nH\bigr).
\]

Indeed, by \SGARefLink{sga3:viii:theorem:6:4}{Exposé VIII, 6.4}, condition \SGARefLink{sga3:xi:proposition:6:1:item:a}{(a)} implies that \({}_nH\) is
finite and flat, and hence a fortiori ``essentially free'' over \(S\).
Thus the right-hand side of \SGARefLink{sga3:xi:display:6:1:cross}{\((x)\)} is representable by a closed
subprescheme \(T_n\) of \(S\).\footnote{The printed French text has
``\(T_n\) of \(T\)'' here; the next sentence makes the \(T_n\) closed
subpreschemes of \(S\), and \(T\) is defined only afterward as their
stable value.}  Order \(E\) by divisibility.  The
\(T_n\) then form a decreasing family of closed subpreschemes of
\(S\); consequently, when \(S\) is noetherian---as we may now
suppose---this family is stationary for sufficiently large \(n\).
Let \(T\) be the stable value of \(T_n\).  We claim that \(T\)
represents the left-hand side of \SGARefLink{sga3:xi:display:6:1:cross}{\((x)\)}, which will complete
the proof.

It is enough to prove the following.  Let \(S'\) be a prescheme over
\(S\) such that
\SGARefTarget{sga3:xi:display:6:1:base-change-condition}
\[
  (Y\cap{}_nH)_{S'}=({}_nH)_{S'}
  \qquad\text{for every \(n\in E\),}
\]
or, equivalently, such that
\(Y_{S'}\supset({}_nH)_{S'}\) for every \(n\in E\).  Then
\(Y_{S'}=H_{S'}\).  This follows at once from \SGARefLink{sga3:ix:lemma:4:4}{Exposé IX, 4.4}:
in view of conditions \SGARefLink{sga3:xi:proposition:6:1:item:a}{(a)} and \SGARefLink{sga3:xi:proposition:6:1:item:b}{(b)}, the family of the
\(({}_nH)_{S'}\) is schematically dense in \(H_{S'}\).

\medskip
\noindent\SGARefTarget{sga3:xi:theorem:6:2}\textbf{Theorem 6.2.} ---
\textit{Let \(S\) be a prescheme, let \(H,E\) be as in
\SGARefLink{sga3:xi:proposition:6:1}{6.1} and satisfy conditions
\SGARefLink{sga3:xi:proposition:6:1:item:a}{(a)} and
\SGARefLink{sga3:xi:proposition:6:1:item:b}{(b)}, let
\(u:H\to G\) be a homomorphism of \(S\)-groups from \(H\) to an
\(S\)-group prescheme locally of finite type over \(S\), and let
\(K\) be a closed subgroup prescheme of \(G\).\footnote{The printed
French text has \(H\) here; the proof immediately uses \(K\subset G\),
as required by the transporter.}  Consider the
subfunctor \(\operatorname{Transp}_G(u,K)\) of \(G\) (cf.\ \SGARefLink{sga3:xi:proposition:2:5}{2.5}).
It is representable by a closed subprescheme of \(G\).  If \(G\) is
noetherian---for example, if \(S\) is noetherian and \(G\) is of
finite type over \(S\)---there exists \(n\in E\) such that}
\SGARefTarget{sga3:xi:display:6:2:transporter-stabilization}
\[
  \operatorname{Transp}_G(u,K)
  =
  \operatorname{Transp}_G\bigl(u|_{{}_nH},K\bigr).
\]
\textit{Suppose, finally, that \(G\) is smooth over \(S\), that
\(K\) is smooth over \(S\) or is of multiplicative type, and that
\(H,E\) satisfy the following condition, stronger than \SGARefLink{sga3:xi:proposition:6:1:item:a}{(a)}:}

\begin{enumerate}
\item[\(a'\)] \SGARefTarget{sga3:xi:theorem:6:2:item:a-prime}%
\textit{for every \(n\in E\), the subgroup \({}_nH\) of \(H\) is
of multiplicative type.}
\end{enumerate}

\textit{Then \(\operatorname{Transp}_G(u,K)\) is smooth over \(S\).}

Consider the \(G\)-group
\SGARefTarget{sga3:xi:display:6:2:hg}
\[
  H_G=H\times_SG,
\]
which plainly satisfies the conditions of
\SGARefLink{sga3:xi:proposition:6:1}{6.1}, with \(S\) replaced by
\(G\) and \(H\) by \(H_G\).  Let \(Y\) be the closed subprescheme of
\(H_G\) that is the inverse image of \(K\subset G\) under the map
\SGARefTarget{sga3:xi:display:6:2:conjugation-map}
\[
  H_G=H\times_SG\longrightarrow G,\qquad
  (h,g)\longmapsto\operatorname{int}(g)\cdot u(h).
\]
Then
\SGARefTarget{sga3:xi:display:6:2:product-identification}
\[
  \operatorname{Transp}_G(u,K)
  =
  \prod_{H_G/G}(Y/H_G)
\]
(compare \SGARefLink{sga3:viii:examples:6:5:item:e:01}{Exposé VIII, Example 6.5(e)}).  The first assertions
therefore follow from \SGARefLink{sga3:xi:proposition:6:1}{6.1}.
Moreover, for every quasi-compact open \(U\subset G\), there exists
\(n\in E\) for which the restrictions to \(U\) of the two functors
\SGARefTarget{sga3:xi:display:6:2:local-stabilization}
\[
  \operatorname{Transp}_G(u,K)
  \qquad\text{and}\qquad
  \operatorname{Transp}_G\bigl(u|_{{}_nH},K\bigr)
\]
coincide.  To verify the last assertion of
\SGARefLink{sga3:xi:theorem:6:2}{6.2}, we may therefore replace
\(H\) by some \({}_nH\), after which it is enough to apply \SGARefLink{sga3:xi:proposition:2:5}{2.5},
since \({}_nH\) is assumed to be of multiplicative type over \(S\).

\medskip
\noindent\SGARefTarget{sga3:xi:remark:6:3}\textbf{Remark 6.3.} ---
The preceding proof uses only the very elementary result of \SGARefLink{sga3:viii:theorem:6:4}{Exposé
VIII, 6.4}, together, for its last part, with \SGARefLink{sga3:xi:proposition:2:5}{2.5}; when \(K\) is
smooth over \(S\), the latter is the infinitesimal result of \SGARefLink{sga3:ix:theorem:3:6}{Exposé
IX, 3.6}, and hence the vanishing of the cohomology of
multiplicative-type groups.  In the most important cases (cf.\
\SGARefLink{sga3:xi:unit:6:7}{6.7}), the integers \(n\in E\) may
even be taken prime to the residue characteristics of \(S\), that
is, invertible in \(\mathscr O_S\).  The \({}_nH\) are then finite
étale over \(S\), and the cohomological result invoked is practically
trivial.  Thus, in that situation,
\SGARefLink{sga3:xi:theorem:6:2}{6.2} is independent of the theory of
multiplicative-type groups.

\medskip
\noindent\SGARefTarget{sga3:xi:unit:6:4}\textbf{6.4.} ---
As usual, one sees that \SGARefLink{sga3:xi:proposition:6:1}{6.1}
extends to the following situation.  Let \(S'\) be an \(S\)-prescheme
(not necessarily locally noetherian), put \(H'=H_{S'}\), and let
\(Y'\) be a closed subprescheme of \(H'\), provided that
\(Y'\to H'\) is of finite presentation, that is, that the ideal
defining \(Y'\) is of finite type.  Then
\(\prod_{H'/S'}(Y'/H')\) is representable by a closed subprescheme
\(T'\) of \(S'\), the morphism \(T'\to S'\) is of finite
presentation, and, if \(S'\) is quasi-compact, there exists \(n\in E\)
for which the analogue of \SGARefLink{sga3:xi:display:6:1:cross}{\((x)\)} holds.  It follows as well
that the first assertion of \SGARefLink{sga3:xi:theorem:6:2}{6.2}
remains valid without assuming that \(G\) is locally of finite type
over \(S\), provided that the immersion \(K\to G\) is of finite
presentation.

\medskip
\noindent\SGARefTarget{sga3:xi:unit:6:5}\textbf{6.5.} ---
As announced in \SGARefLink{sga3:xi:remark:5:6}{5.6}, Theorem
\SGARefLink{sga3:xi:theorem:6:2}{6.2} makes it possible to extend to
group preschemes satisfying \SGARefLink{sga3:xi:theorem:6:2:item:a-prime}{\((a')\)} and \SGARefLink{sga3:xi:proposition:6:1:item:b}{(b)} above certain results
proved by another method, under more restrictive hypotheses, for
groups of multiplicative type.  This applies to \SGARefLink{sga3:xi:corollary:5:2}{5.2}, the beginning
of \SGARefLink{sga3:xi:corollary:5:3}{5.3}, and \SGARefLink{sga3:xi:corollary:5:4}{5.4} and their corresponding ``bis'' variants, and also
to \SGARefLink{sga3:xi:proposition:5:9}{5.9} and \SGARefLink{sga3:xi:corollary:5:10}{5.10}.  It also applies to the results of
\SGARefLink{sga3:ix:section:5}{Exposé IX, Sections 5} and 6, with the exception of \SGARefLink{sga3:ix:theorem:6:8}{Exposé IX, 6.8}.
The latter assertion is already false for a homomorphism of abelian
schemes \(u:H\to G\) over the spectrum \(S\) of a discrete valuation
ring of residue characteristic \(p>0\): it can happen that
\(\operatorname{Ker}u\) has the trivial group as its generic fibre
and, as its special fibre, a radicial group that is not reduced to
the trivial group.

\medskip
\noindent\SGARefTarget{sga3:xi:unit:6:6}\textbf{6.6.} ---
We have just given an example of a rigidity property of
multiplicative-type groups that is not shared by abelian schemes.
Another, extremely important, example is that the existence theorem
for infinitesimal extensions of homomorphisms in \SGARefLink{sga3:ix:theorem:3:6}{Exposé IX, 3.6} no
longer holds when \(H\) is an abelian scheme.  Thus a nonzero abelian
scheme over a field admits nontrivial infinitesimal variations, in
contrast to a multiplicative-type group.  This is the infinitesimal
aspect of the fact that there is a theory of moduli---one that,
moreover, is far from complete---for abelian varieties, whereas the
theory of moduli of multiplicative-type groups is empty.

Another, ``global,'' aspect of this infinitesimal difference is the
following.  Let \(H\) be an abelian scheme over a locally noetherian
prescheme \(S\), and let \(G\) be a commutative \(S\)-group prescheme
locally of finite type over \(S\).  One can show that
\(\operatorname{Hom}_{S\text{-}\mathrm{gr}}(H,G)\) is representable
by a prescheme locally of finite type over \(S\).  Unlike the
corresponding prescheme when \(H\) is of multiplicative type,
however, it is not étale over \(S\), but only unramified.  Thus, for
example, if \(S\) is the spectrum of a complete discrete valuation
ring and \(H,G\) are abelian schemes over \(S\), there may exist
homomorphisms \(H_0\to G_0\) on the special fibres that do not arise
``by specialization'' from a homomorphism on the generic fibres.

\medskip
\noindent\SGARefTarget{sga3:xi:unit:6:7}\textbf{6.7.} ---
Theorem \SGARefLink{sga3:xi:theorem:6:2}{6.2} applies whenever \(H\)
is an abelian scheme over \(S\), or, more generally, an extension of
such a scheme by a torus.  Indeed, the question is local on \(S\), so
one may suppose that there is a prime number \(\ell\) prime to the
residue characteristics of \(S\).  Taking \(E\) to be the set of
powers of \(\ell\) then suffices to satisfy conditions \SGARefLink{sga3:xi:theorem:6:2:item:a-prime}{\((a')\)} and
\SGARefLink{sga3:xi:proposition:6:1:item:b}{(b)}.

An argument analogous to that of
\SGARefLink{sga3:xi:proposition:6:1}{6.1} gives \SGARefLink{sga3:xi:theorem:6:8}{the following result}.

\medskip
\noindent\SGARefTarget{sga3:xi:theorem:6:8}\textbf{Theorem 6.8.} ---
\textit{Let \(X\) be a prescheme smooth over \(S\), with nonempty
connected geometric fibres.  For every closed subprescheme \(Y\) of
\(X\), the functor \(\prod_{X/S}(Y/X)\) is representable by a closed
subprescheme \(T\) of \(S\).  If \(Y\) is of finite presentation
over \(X\), then \(T\) is of finite presentation over \(S\).}

Since \(f:X\to S\) is faithfully flat and locally of finite
presentation, it is a covering for the fpqc topology.  On the other
hand, \(T=\prod_{X/S}(Y/X)\) is plainly a subsheaf of \(S\) for the
fpqc topology.  Hence both the representability of \(T\) by a closed
subprescheme of \(S\) and the question whether \(T\) is of finite
presentation over \(S\) are local on \(S\) for the fpqc topology.
After the base change \(S'\to S\) with \(S'=X\), we are reduced to
the case in which \(X\) admits a section \(e\) over \(S\).  We may
also suppose that \(S\) is affine, and hence in particular
quasi-compact.  We then have:

\medskip
\noindent\SGARefTarget{sga3:xi:corollary:6:9}\textbf{Corollary 6.9.}
---
\textit{Under the hypotheses of
\SGARefLink{sga3:xi:theorem:6:8}{6.8}, suppose that \(S\) is
quasi-compact and that \(X\) admits a section \(e\) over \(S\).  For
every integer \(n\geq0\), let \(X_n\) be the order-\(n\)
infinitesimal neighbourhood of the section \(e\) in \(X\).  Suppose
that \(Y\) is of finite presentation over \(X\).  Then there exists
an integer \(n\geq0\) such that}
\SGARefTarget{sga3:xi:display:6:9:stabilization}
\[
  \prod_{X/S}(Y/X)
  =
  \prod_{X_n/S}(Y_n/X_n),
  \qquad Y_n=Y\cap X_n.
\]

This corollary implies
\SGARefLink{sga3:xi:theorem:6:8}{6.8} when \(Y\) is of finite
presentation over \(X\), by \SGARefLink{sga3:viii:theorem:6:4}{Exposé VIII, 6.4}.  Indeed, since \(X\)
is smooth over \(S\), \(X_n\) is finite and locally free over \(S\),
and a fortiori ``essentially free'' over \(S\) in the sense of
\SGARefLink{sga3:viii:definition:6:1}{Exposé VIII, 6.1}.  Thus \(\prod_{X_n/S}(Y_n/X_n)\) is representable
by a closed subprescheme of \(S\).  Furthermore, either the proof of
the cited result or reduction to the noetherian case shows at once
that this closed subprescheme is of finite presentation over \(S\).

Let us first prove \SGARefLink{sga3:xi:corollary:6:9}{6.9}, and hence
\SGARefLink{sga3:xi:theorem:6:8}{6.8}, when \(S\) is noetherian.  Put
\SGARefTarget{sga3:xi:display:6:9:tn}
\[
  T_n=\prod_{X_n/S}(Y_n/X_n).
\]
The \(T_n\) form a decreasing sequence of closed subpreschemes of
\(S\), and this sequence is stationary because \(S\) is noetherian.
Let
\SGARefTarget{sga3:xi:display:6:9:r}
\[
  R=\bigcap_{n\geq0}\prod_{X_n/S}(Y_n/X_n)
\]
be their common value for large \(n\).  Plainly \(T\subset R\), and
it remains to show that \(R\subset T\).  After the base change
\(R\to S\), we are reduced to the case \(R=S\), that is,
\(Y_n=X_n\) for every \(n\), or equivalently
\(Y\supset X_n\) for every \(n\).  We must then prove \(T=S\), that
is, \(Y=X\).

The inclusions \(Y\supset X_n\) for all \(n\), together with the fact
that \(X\) is locally noetherian, imply that in a neighbourhood of
each point of \(e(S)\), \(Y\) is an induced open subprescheme of
\(X\).  Hence there is an induced open subprescheme \(U\) of \(X\)
such that \(e(S)\subset U\subset Y\).  By \SGARefLink{sga3:ix:corollary:4:3}{Exposé IX, 4.3}, since the
fibres of \(X/S\) are integral, \(U\) is schematically dense in
\(X\).  As \(Y\) is a closed subprescheme containing \(U\), it
follows that \(Y=X\).  This proves
\SGARefLink{sga3:xi:corollary:6:9}{6.9}, and hence
\SGARefLink{sga3:xi:theorem:6:8}{6.8}, in the noetherian case.

The general case is reduced to the preceding one.  For every
\(s\in S\), there is an affine open neighbourhood \(U\) of \(s\)
and an affine open neighbourhood \(V\) of \(e(s)\) such that
\(f(V)\subset U\).  Then \(f(V)\) is an open neighbourhood of \(s\)
contained in \(U\).  Choose an affine open neighbourhood \(S'\) of
\(s\) contained in \(e^{-1}(V)\cap f(V)\), and put
\(X'=V\cap f^{-1}(S')\).  Then \(X'\) and \(S'\) are affine open
subpreschemes of \(X\) and \(S\), respectively, and \(X'/S'\)
admits a section.  By the local nature of
\SGARefLink{sga3:xi:theorem:6:8}{6.8} and
\SGARefLink{sga3:xi:corollary:6:9}{6.9}, we may suppose that
\(S=S'\).

We claim that
\SGARefTarget{sga3:xi:display:6:9:open-reduction}
\[
  \prod_{X/S}(Y/X)=\prod_{X'/S'}(Y'/X'),
  \qquad Y'=Y\cap X'.
\]
Indeed, by \SGARefLink{sga3:ix:corollary:4:6}{Exposé IX, 4.6}, \(X'\) is schematically dense in \(X\),
at least when \(X\) is quasi-separated over \(S\), so that \(X'\)
is retrocompact in \(X\).  In fact, one can readily prove that
\SGARefLink{sga3:ix:corollary:4:6}{Exposé IX, 4.6} remains valid without the retrocompactness
hypothesis.  Likewise, after every base change \(S_1\to S\),
\(X'\times_SS_1\) is schematically dense in \(X\times_SS_1\), and
the claimed equality follows immediately.  We are thus reduced to
the case \(X=X'\), so we may suppose that both \(S\) and \(X\) are
affine.

Now let \(X=\operatorname{Spec}(B)\), and let \(J\) be the ideal of
\(B\) defining \(Y\).  The ideal \(J\) is the filtered colimit of
its finite-type subideals.  Consequently, \(Y\) is the intersection
of closed subpreschemes \(Y_i\) of \(X\) that are of finite
presentation over \(S\), and
\SGARefTarget{sga3:xi:display:6:9:intersection}
\[
  \prod_{X/S}(Y/X)
  =
  \bigcap_i\prod_{X/S}(Y_i/X).
\]
To prove \SGARefLink{sga3:xi:theorem:6:8}{6.8}, it is therefore
enough to treat the case in which \(Y\) is of finite presentation
over \(X\).  It remains to prove
\SGARefLink{sga3:xi:corollary:6:9}{6.9} with \(S\) and \(X\)
affine.  But \(X\) and \(Y\) over \(S\) then arise by base change
\(S\to S_0\) from an analogous situation \(X_0,Y_0\) over a
noetherian prescheme \(S_0\).  This reduces us to the already treated
case in which \(S\) is noetherian, and completes the proof of
\SGARefLink{sga3:xi:theorem:6:8}{6.8} and
\SGARefLink{sga3:xi:corollary:6:9}{6.9}.

\medskip
\noindent\SGARefTarget{sga3:xi:corollary:6:10}\textbf{Corollary 6.10.}
---
\textit{Let \(X\) be an \(S\)-group prescheme, smooth and of finite
presentation, with connected fibres.  Let \(Y\) be a group
prescheme of finite presentation over \(S\), and let \(i:Y\to X\)
be a monomorphism of \(S\)-group preschemes, so that \(Y\) is a
subgroup of \(X\).  Then \(\prod_{X/S}(Y/X)\) is representable by a
closed subprescheme of \(S\) of finite presentation.  If \(S\) is
quasi-compact, let \(X_n\) denote, for every \(n\geq0\), the
order-\(n\) infinitesimal neighbourhood of the unit section of
\(X\), and put \(Y_n=X_n\cap Y\).  Then, for all sufficiently large
\(n\),}
\SGARefTarget{sga3:xi:display:6:10:stabilization}
\[
  \prod_{X/S}(Y/X)=\prod_{X_n/S}(Y_n/X_n).
\]

The proof is essentially that of
\SGARefLink{sga3:xi:corollary:6:9}{6.9}.  The unit sections of \(X\)
and \(Y\) induce bijective immersions \(S\to X_n\) and
\(S\to Y_n\), and hence isomorphisms
\SGARefTarget{sga3:xi:display:6:10:reductions}
\[
  S_{\mathrm{red}}\xrightarrow{\sim}(X_n)_{\mathrm{red}},
  \qquad
  S_{\mathrm{red}}\xrightarrow{\sim}(Y_n)_{\mathrm{red}}.
\]
It follows that the injection \(Y_n\to X_n\) is proper.  Since it is
a monomorphism of finite presentation, it is therefore a closed
immersion.  \SGARefLink{sga3:viii:theorem:6:4}{Exposé VIII, 6.4}, already used above, now shows that
\(\prod_{X_n/S}(Y_n/X_n)\) is representable by a closed
subprescheme of \(S\) of finite presentation over \(S\).  It remains
to prove the last assertion of
\SGARefLink{sga3:xi:corollary:6:10}{6.10} under the additional
assumption that \(S\) is affine.

We reduce once more to the case in which \(S\) is noetherian.  With
the notation of the proof of
\SGARefLink{sga3:xi:corollary:6:9}{6.9}, it is enough to prove
\(R=T\), or, equivalently, that \(Y\supset X_n\) for every \(n\)
implies \(Y=X\).  This hypothesis implies that \(i:Y\to X\) is
étale at the points of the unit section of \(Y\) over \(S\), and
therefore that \(Y\) is smooth over \(S\) at those points.  Hence the
open subprescheme \(Y'\) consisting of the points where \(Y\) is
smooth over \(S\) is an induced open subgroup of \(Y\).  By \SGARefLink{sga3:x:lemma:3:5}{Exposé
X, 3.5}, the morphism \(Y'\to X\) is an étale monomorphism and hence
an open immersion.  Since the fibres of \(X\) are connected and
every open subgroup of an algebraic group is also closed, this open
immersion is surjective, and therefore an isomorphism.  Thus
\(Y'=X\), and a fortiori \(Y=X\), completing the proof of
\SGARefLink{sga3:xi:corollary:6:10}{6.10}.

Arguing as in \SGARefLink{sga3:viii:examples:6:5}{Exposé VIII, 6.5}, we deduce:

\medskip
\noindent\SGARefTarget{sga3:xi:corollary:6:11}\textbf{Corollary 6.11.}
---
\textit{Let \(G\) be an \(S\)-group prescheme locally of finite type
and quasi-separated over \(S\), let \(H\) be a group prescheme
smooth and of finite presentation over \(S\), with connected
fibres, and let \(i:H\to G\) be a monomorphism of \(S\)-groups, so
that \(H\) is a subgroup of \(G\).  Then:}

\begin{enumerate}
\item[(a)] \SGARefTarget{sga3:xi:corollary:6:11:item:a}%
\textit{\(\operatorname{Centr}_G(H)\) and
\(\operatorname{Norm}_G(H)\) are representable by closed
subpreschemes of \(G\) of finite presentation over \(G\).  Likewise,
for every monomorphism \(j:K\to G\) of finite presentation between
\(S\)-group preschemes, \(\operatorname{Transp}_G(H,K)\) is
representable by a closed subprescheme of \(G\) of finite
presentation over \(G\).}

\item[(b)] \SGARefTarget{sga3:xi:corollary:6:11:item:b}%
\textit{If \(G\) is quasi-compact, then in each of the cases in
\SGARefLink{sga3:xi:corollary:6:11:item:a}{(a)} there exists an integer \(n\geq0\) such that, if \(H_n\)
denotes the order-\(n\) infinitesimal neighbourhood of the unit
section of \(H\) (and \(K_n\) is defined similarly in the transporter
case),}
\SGARefTarget{sga3:xi:display:6:11:stabilization}
\[
  \begin{aligned}
  \operatorname{Centr}_G(H)
    &=\operatorname{Centr}_G(H_n),\\
  \operatorname{Norm}_G(H)
    &=\operatorname{Norm}_G(H_n),\\
  \operatorname{Transp}_G(H,K)
    &=\operatorname{Transp}_G(H_n,K)
     =\operatorname{Transp}_G(H_n,K_n).
  \end{aligned}
\]
\end{enumerate}

Apply \SGARefLink{sga3:xi:corollary:6:10}{6.10} over the base
prescheme \(G\) to the group prescheme
\SGARefTarget{sga3:xi:display:6:11:hg}
\[
  X=H_G=H\times_SG.
\]
For the centralizer, take as \(Y\) the inverse image of the diagonal
subgroup of \((G\times_SG)_G\) under a suitable homomorphism of
\(G\)-groups from \(X\) to \((G\times_SG)_G\).  For the transporter,
take the inverse image of \(K_G\) under a suitable homomorphism of
\(G\)-groups from \(X\) to \(G_G\).  The normalizer case reduces to
the transporter case by taking \(K=H\).  The hypothesis on \(G\)
ensures that \(H\to G\) is of finite presentation, and hence that
\(Y\to X\) is of finite presentation.  In the centralizer case, the
same hypothesis ensures that the diagonal subgroup of
\(G\times_SG\) is of finite presentation over \(G\times_SG\);
again it follows that \(Y\) is of finite presentation over \(X\).

\medskip
\noindent\SGARefTarget{sga3:xi:remark:6:12}\textbf{Remark 6.12.} ---
One can prove, using rather delicate results of EGA
VI,\footnote{Cf. also J.\,P. Murre, ``Representation of unramified
functors. Applications,'' \emph{Séminaire Bourbaki}, no.\ 294
(May 1965), Theorem 3, p.\ 13.} that if \(X\) is a prescheme flat
and of finite presentation over \(S\), and is either proper over
\(S\) or has nonempty connected geometric fibres, then for every
closed subprescheme \(Y\) of \(X\) of finite presentation over
\(S\), the functor \(\prod_{X/S}(Y/X)\) is representable by a closed
subprescheme of \(S\) of finite presentation over \(S\).  Likewise,
if \(X\) is an \(S\)-group prescheme flat and of finite presentation
with connected fibres, and \(i:Y\to X\) is a monomorphism of
\(S\)-groups, with \(Y\) an \(S\)-group prescheme of finite
presentation, then \(\prod_{X/S}(Y/X)\) is representable by a closed
subprescheme of \(S\) of finite presentation over \(S\).  In
particular, \SGARefLink{sga3:xi:corollary:6:11:item:a}{6.11(a)}
remains valid if the hypothesis ``\(H\) smooth over \(S\)'' is
replaced by ``\(H\) flat over \(S\).''

